\section{High cutoffs, actual bilinear factors and a signed covariance test}\label{sb:core}
Retain the original Eisenstein family, fixed finite-order $\nu$, fixed $S$,
common $W\in C_c^\infty((0,\infty))$ supported in $[1,2]$, and every raw
shared-prime zero. Put
\[
 C_U(m)=\sum_{\substack{abc=m\\Na>U,\,Nb>U}}\mu(a)\mu(b),\qquad
 r_U(m)=\#\{(a,b):ab=m,\ Na,Nb>U\}.
\]
All factors are ideals, squarefree and pairwise coprime; on the active
annulus $A_U=C_U-\mu$. A coefficient identity is different from a row-energy bound.

\begin{proposition}[The high-cutoff window]\label{sb:window}
If $D>1$ and $2D/7\le U^2\le D$, then for every surviving squarefree
column $D\le Nm\le2D$,
\[
 C_U(m)=\mu(m)r_U(m),\qquad
 A_U(m)=\mu(m)(r_U(m)-1),\qquad |A_U(m)|\ge1.       
\]
For each fixed nonzero $W$, the normalized squared coefficient norm satisfies
\begin{equation}
 D^{-1}\sum_{\substack{m\ {\rm sf}\\(m,S)=1}}
 |A_U(m)W(Nm/D)|^2\gg_{W,S}1                     \label{sb:coefficient-lower}
\end{equation}
for large $D$, uniformly in this window. This is not a lower bound for row energy.
\end{proposition}
\begin{proof}
The least norm of a nonunit ideal outside $S$ is at least $7$. A long-long
term with $c\ne1$ would have $Nm>7U^2\ge2D$. Thus $ab=m$, including both
endpoints. The fixed-point-free swap $(a,b)\mapsto(b,a)$ makes $r_U$ even,
so $A_U$ is an odd nonzero integer. The last claim follows by counting
squarefree ideals on a fixed interval where $|W|$ is bounded below; the
constant and error are proved in the extended proof.
\end{proof}

The exact remaining polynomial is genuinely bilinear, with $(a,b)=1$,
$U<Na,Nb<2D/U$, $D\le Na\,Nb\le2D$, and
$x_{ab}(u)=x_a(u)x_b(u)$, including zeros. Both factors see the same row
$u=\varepsilon\alpha\beta^2$. Cauchy followed by the accepted powered BGL
bound gives only
\begin{equation}
 Q_{B_X}(C_U,C_U)\ll_{\nu,S,\epsilon}
 (DX)^\epsilon\|W\|_\infty^2
 \{\sqrt D\,X^2+DX+D^{5/6}X^{5/3}\}.              \label{sb:factor-bound}
\end{equation}
At $X\asymp D^{1/4}$ the exponents are $1,5/4,5/4$; the earlier regrouped
bound has $1/2,5/4,13/12$ and is stronger. No $BT_\sigma$ saving is obtained.
Multiplying two marginal operator norms does not supply a coupled fourth moment.

\paragraph{A finite test in the original family.}
There is no automatic nonnegative signed covariance, even when
$C_U(m)\mu(m)\ge0$. Take $\nu=1$, $S$ the primes above $2,3$, $D=400$,
$U=17$, and all $72$ literal rows of $B_7$. Choose the common nonnegative
smooth profile constructed in the extended proof, with values $10$ at $439/400$, $1$ at
$475/400$, and zero at every other integer norm divided by $400$.
There are exactly four active squarefree ideal columns, of norms
$439,439,475,475$. In that order,
\[
 A=(1,1,1,1),\quad C=(0,0,2,2),\quad
 \mu=(-1,-1,1,1),\qquad
 \Re G_{B_7}=
 \begin{pmatrix}72&0&0&18\\0&72&18&0\\0&18&72&0\\18&0&0&72\end{pmatrix}.
\]
Exact raw-symbol enumeration, independently reproducible from the ideal
and residue-field definitions, gives
\begin{equation}
 Q(C,\mu)=-\frac{27}{25},\qquad
 Q(C,C)=\frac{36}{25},\quad
 Q(\mu,\mu)=\frac{864}{25},\quad
 Q(A,A)=\frac{954}{25}.                          \label{sb:covariance-example}
\end{equation}
The real Gram matrix is positive definite, with eigenvalues $54,90$;
its positivity does not force a cross covariance to be positive.
This fixed finite example refutes that shortcut, not the asymptotic
$BT_\sigma$ target or the original signed mean-square estimate.

\paragraph{The two-cutoff contour remains coupled.}
For $T_0=S\cup\operatorname{supp}(u)$, $\psi=\psi_u$ and independently
primitive $\psi_2=\operatorname{prim}(\psi^2)$, put $q_p=\psi(p)(Np)^{-s}$,
$\alpha_p=(Np)^{-z}$, $\beta_p=(Np)^{-w}$ outside $T_0$. The actual
three-placement Euler factor is $1+q_p-q_p\alpha_p-q_p\beta_p$. Its exact
factorization is
\begin{equation}
 \mathcal F(s,z,w)=
 \frac{L^{T_0}(s,\psi)}{L^{T_0}(s+z,\psi)L^{T_0}(s+w,\psi)}
 \prod_{p\notin T_0}\left(1-
 \frac{q_p^2(1-\alpha_p)(1-\beta_p)}
 {(1-q_p\alpha_p)(1-q_p\beta_p)}\right).           \label{sb:perron-factor}
\end{equation}
Exact double Perron inversion is valid on positive auxiliary lines
$c,d>1-\Re s$ for $1/2<\Re s<1$, retaining the original principal pole
exclusion. These safe lines reproduce the previous $U^{2-4\delta}$
energy loss at $\Re s=1/2+\delta$, with the same common-profile derivative
and fixed-margin dependencies. The correction product alone does not
justify moving those lines through reciprocal-$L$ poles. In fact
$\mathcal F(s,z,0)=1/L^{T_0}(s+z,\psi)$ as a meromorphic identity.
No zero-free region, reciprocal moment, analytic saving, new Lean
formalization or GRH proof is supplied by this construction.
